\section{New Delhi Case Study and Experimental Platform}
\label{sec:new_delhi_platform}

We instantiate the framework on a single New Delhi platform so that the
controlled and observed-data studies share one regulatory sensor geometry, one
PM\(_{2.5}\) and wind record, and one set of proxy inventories;
Section~\ref{sec:evaluation} then varies wind, noise, backgrounds, source
ambiguity, activity bases, and forward-model error within it.

\subsection{Government Observations and Sensor Network}
\label{subsec:new_delhi_observations}

We use hourly government observations \citep{cpcb,cpcb_portal} (columns
\texttt{monitor\_id}, \texttt{timestamp\_round}, \texttt{AT}, \texttt{RH},
\texttt{WD}, \texttt{WS}, \texttt{pm10}, \texttt{pm25}) spanning 21,960 hourly
timestamps from 1 May 2018 through 31 October 2020 in Indian Standard Time;
approximately 74.9\% of wind-direction, 72.8\% of wind-speed, and 90.5\% of
PM\(_{2.5}\) entries are observed, and a valid-value mask (\texttt{True} where the
record is present) is retained alongside station coordinates and timestamps.  The
two Pusa monitors are not treated as coincident independent sensors:
\texttt{Pusa\_IMD} and \texttt{Pusa\_DPCC} are averaged by timestamp into
\texttt{Pusa\_averaged} (mean coordinates, per-variable hourly averaging), giving
the final 32-sensor layout.  PM\(_{2.5}\) is never imputed: its flattened
valid-value mask defines \(M_{\mathcal O}\) in
Equation~\ref{eq:observation_selection}, and the same ordered sensor--time rows
are retained in \(Y\), \(H_\Phi^{\mathrm{lag}}\), \(Q\), and their metadata.

\subsection{Wind Preparation}
\label{subsec:new_delhi_wind}

The observed \(WD/WS\) values are converted using
Equation~\ref{eq:wind_direction_conversion} and completed on the response grid by
the kernel coordinate-query imputer of
Section~\ref{subsec:wind_field_estimation}; a learned imputer was trained and
evaluated but not adopted (Section~\ref{subsec:results_wind_imputation}).  The
saved wind product retains raw station values, masks, station coordinates,
held-out validation splits, checkpoint metadata, and the gridded query field used
by the response operator.

\subsection{Inventories and Proxy Temporal Activities}
\label{subsec:new_delhi_inventories}

The study uses four source groups---brick kilns, industries, population density,
and traffic (Table~\ref{tab:new_delhi_inventories})---each with one spatial proxy
map and its own admissible set of temporal-basis components.  Brick-kiln and
industry proxies follow the regional emissions inventory
\citep{GUTTIKUNDA2013101}, population density uses the gridded population product
\citep{ColumbiaUniversity2018}, and the traffic map is derived from road map data
\citep{google_maps}.  Every \(80\times80\) map is cropped to the New Delhi study
window (rows \(21{:}61\), columns \(16{:}56\)) covering the regulatory network,
giving a \(40\times40\) map, then divided by its own cropped 99th percentile.
Traffic is modeled as a single road-network source whose nearest-slot maps at
00:00, 06:00, 12:00, and 18:00 define traffic-specific temporal-basis components
under the shared dictionary \(\Phi\) and the fixed-zero mask \(\mathcal F_0\), so
one road map carries a time-varying activity (assuming the congestion
\emph{spatial} pattern is time-invariant and only its amplitude varies).
Brick-kiln activity uses deterministic alternating 12-hour blocks, industry a
continuous operating fraction with higher daytime (07:00--19:00) activity, and
population a baseline with smooth cooking peaks---deterministic study assumptions,
not observed emissions truth.  Because each map uses independent proxy
normalization, fitted coefficients and contributions are in normalized proxy
units, not physical emission totals or shares; no learned free residual source
map is used, and unexplained broad signal belongs to the background model or
residual diagnostics.

\begin{table}[t]
\centering
\small
\setlength{\tabcolsep}{4pt}
\caption{New Delhi source groups, spatial proxies, and temporal-basis components.}
\label{tab:new_delhi_inventories}
\begin{tabularx}{\columnwidth}{l
  >{\raggedright\arraybackslash}X
  >{\raggedright\arraybackslash}X}
\toprule
Source group & Spatial proxy & Temporal-basis components \\
\midrule
Brick kilns & regional kiln map & alternating 12-hour blocks \\
Industries & regional industry map & continuous + daytime (07--19) \\
Population density & gridded population & peaks 07:00, 13:00, 19:00 \\
Traffic & road-network map & diurnal slots 00/06/12/18 \\
\bottomrule
\end{tabularx}
\end{table}

\subsection{Forward Models, Sensor Observations, and Backgrounds}
\label{subsec:new_delhi_forward_models}

Both forward paths use the New Delhi-derived \(40\times40\) domain and regulatory
sensor coordinates, with the open-boundary Gaussian puff operator of
Section~\ref{subsec:plume_response} evaluated directly at sensors.  The same
operator is the matched generator for controlled recovery experiments
(Experiments~1--4) and the inference operator for observed PM\(_{2.5}\); only the
initial-condition baseline differs by mode---a zero-source, zero-initial-state
baseline for controlled operator-matched runs, and a subtracted zero-source
transported kriged-initial-condition baseline for observed runs.  An auxiliary
edge-hold advection--diffusion simulator provides a \emph{structural}
forward-model mismatch (data generated by a distinct transport operator, fit with
the puff response) that also drives the residual-adequacy test of
Section~\ref{subsec:model_adequacy} (Experiment~5).  Normal backgrounds follow
Section~\ref{subsec:background_construction}, with the base controlled background
the rank-four member; the kriged-baseline construction and simulator settings are
given in Appendix~\ref{app:platform_details}.

\subsection{Controlled and Observed-Data Study Modes}
\label{subsec:new_delhi_study_modes}

The controlled mode preserves New Delhi sensor geometry and transport while
assigning known nonnegative source--basis coefficients, using synthetic or real
gridded wind, controlled observation noise, source splits, sensor-layout variants,
and optional operator mismatch.  Where the study is explicitly about inventory
geometry (Experiments~2 and~6) the real New Delhi inventory maps are used; the
remaining controlled experiments use compact synthetic source cells for
tractability, since their conclusions concern conditioning, coherence, and
recovery rather than specific inventory footprints.  Because its coefficients are
known, it supports recovery and interval-coverage metrics.  The observed-data mode
combines real PM\(_{2.5}\), its observation mask, the gridded kernel-imputed wind
field, the named normalized inventories, and the declared proxy temporal bases; it
has no ground-truth source activities and is run over four consecutive one-week
windows, each fitted and reported independently, reporting residuals, response
geometry, uncertainty, weak and ambiguous coefficients, normalized proxy
contributions, and recommended report groups rather than synthetic recovery error.
